\chapter{门罗币环机密交易 (RingCT) (Monero Ring Confidential Transactions (RingCT))}
\label{chapter:transactions}

在第 \ref{chapter:addresses} 章和第 \ref{chapter:pedersen-commitments} 章中，我们逐步建立了门罗币交易的多个方面。到目前为止，一个简单的单 input、单 output 的交易——从某个未知作者到某个未知接收者——听起来像这样：

"我的交易使用交易公钥 $r G$。我将花费旧的 output $X$（注意它有一个隐藏的金额 $A_X$，在 $C_X$ 中做出承诺）。我将给它一个伪 output 承诺 $C'_X$。我将创建一个 output $Y$，可由一次性地址 $K^o_Y$ 花费。它有一个隐藏的金额 $A_Y$ 在 $C_Y$ 中做出承诺，为接收者加密，并通过 Bulletproofs 风格的范围证明证明在范围内。请注意 $C'_X - C_Y = 0$。"

还有一些问题。作者是否真的拥有 $X$？伪 output 承诺 $C'_X$ 是否真的对应于 $C_X$，使得 $A_X = A'_X = A_Y$？是否有人篡改了交易，也许将 output 指向了原始作者未打算的接收者？\\

如第 \ref{sec:one-time-addresses} 节所述，我们可以通过使用一次性地址签名消息来证明对 output 的所有权（谁拥有该地址的密钥就拥有该 output）。我们还可以通过证明对零承诺的私钥的知识来证明它与伪 output 承诺有相同的金额（$C_X - C'_X = z_X G$）。此外，如果该消息是{所有交易数据}（签名本身除外），那么验证者可以确信一切都是按作者的意图（签名只对原始消息有效）。MLSAG 签名允许我们做到所有这些，同时还将实际花费的 output 混入区块链中的其他 output 中，这样观察者就无法确定正在花费的是哪一个。



\section{交易类型 (Transaction types)}
\label{sec:transaction_types}

门罗币\marginnote{src/crypto- note\_core/ cryptonote\_ tx\_utils.cpp {\tt construct\_ tx\_with\_ tx\_key()}}是一种持续发展中的加密货币。交易结构、协议和密码方案总是会随着新的创新、目标或威胁的发现而演变。

在本报告中，我们将注意力集中在{环机密交易 (Ring Confidential Transactions)}，即 {RingCT}，因为它们是在当前版本的门罗币中实现的。RingCT 对所有新的门罗币交易都是强制性的，因此我们不会描述任何已弃用的交易方案，即使它们仍然被部分支持。\footnote{RingCT 于 2017 年 1 月首次实现（协议 v4）。它在 2017 年 9 月（协议 v6）被强制要求用于所有新交易 \cite{ringct-dates}。RingCT 是门罗币交易协议的第 2 版。}我们到目前为止讨论的交易类型，以及将在本章中进一步详细说明的，是 {\tt RCTTypeBulletproof2}。\footnote{在 RingCT 时代有三种已弃用的交易类型：{\tt RCTTypeFull}、{\tt RCTTypeSimple} 和 {\tt RCTTypeBulletproof}。前两种在 RingCT 的第一次迭代中并存，本报告的第一版 \cite{ztm-1} 中有探讨，然后随着 Bulletproofs 的出现（协议 v8）{\tt RCTTypeFull} 被弃用，{\tt RCTTypeSimple} 被升级为 {\tt RCTTypeBulletproof}。{\tt RCTTypeBulletproof2} 由于对 output 承诺的掩码和金额进行加密的小改进（v10）而到来。}

我们在第 \ref{sec:transaction_summary} 节中给出交易的概念性总结。



\section{类型为 {\tt RCTTypeBulletproof2} 的环机密交易 (Ring Confidential Transactions of type {\tt RCTTypeBulletproof2})}
\label{sec:RCTTypeBulletproof2}

目前（协议 v12）所有新交易都必须使用这种交易类型，其中每个 input 是单独签名的。一个 {\tt RCTTypeBulletproof2} 交易的实际示例，包含所有组件，可以在附录 \ref{appendix:RCTTypeBulletproof2} 中查看。


\subsection{金额承诺与交易手续费 (Amount commitments and transaction fees)}
\label{sec:commitments-and-fees}

假设一个交易发送者之前收到了各种 output，金额为 $a_1, ..., a_m$，地址为一次性地址 $K^o_{\pi,1}, ..., K^o_{\pi,m}$，金额承诺为 $C^a_{\pi,1}, ..., C^a_{\pi,m}$。

该发送者知道对应于一次性地址的私钥 $k^o_{\pi,1}, ..., k^o_{\pi,m}$（第 \ref{sec:one-time-addresses} 节）。发送者还知道承诺 $C^a_{\pi,j}$ 中使用的盲因子 $x_j$（第 \ref{sec:pedersen_monero} 节）。

通常，交易 output 的总额{低于}交易 input，以提供手续费来激励矿工将交易纳入区块链。\footnote{在门罗币中，有一个最低基础手续费，按交易权重缩放。它是半强制性的，因为虽然你可以创建包含微小手续费交易的新区块，但大多数门罗币节点不会将此类交易转发给其他节点。结果是很少有矿工会尝试将它们纳入区块。如果交易作者愿意，可以提供高于最低值的矿工手续费。我们在第 \ref{subsec:dynamic-minimum-fee} 节中对此进行了更详细的讨论。}交易手续费金额 $f$ 以明文存储在发送给网络的交易数据中。矿工可以为自己创建一个额外的 output 来收取手续费（见第 \ref{subsec:miner-transaction} 节）。

一笔交易由 input \(a_1, ..., a_m\) 和 output \(b_0, ..., b_{p-1}\) 组成，使得 \(\sum\limits_{j=1}^m a_j - \sum\limits_{t=0}^{p-1} b_t - f = 0\)。\footnote{Output\marginnote{src/crypto- note\_core/ cryptonote\_ tx\_utils.cpp {\tt construct\_ tx\_with\_ tx\_key()}}[-1.9cm] 在被分配索引之前由核心实现随机打乱，因此观察者无法围绕其顺序建立启发式规则。Input 在交易数据中按 key image 排序。}

发送者为 input 金额计算伪 output 承诺 $C'^a_{\pi,1}, ..., C'^a_{\pi,m}$，并为预期的 output 金额 $b_0, ..., b_{p-1}$ 创建承诺。设这些新承诺为 $C^b_0, ..., C^b_{p-1}$。

他知道零承诺 $(C^a_{\pi,1} - C'^a_{\pi,1}),...,(C^a_{\pi,m} - C'^a_{\pi,m})$ 的私钥 $z_1,...,z_m$。

为了让验证者确认交易金额之和为零，手续费金额必须转换为承诺。解决方案是在没有任何盲因子的掩蔽效果下计算手续费 $f$ 的承诺。即 $C(f) = f H$。

现在\marginnote{src/ringct/ rctSigs.cpp verRct- Semantics- Simple()}我们可以证明 input 金额等于 output 金额：\\
\[(\sum_j C'^a_{j} - \sum_t C^b_{t}) - f H = 0\]


\subsection{签名 (Signature)}
\label{full-signature}

发送者从区块链中选择 $m$ 组、每组 $v$ 个额外的不相关一次性地址及其承诺，对应于看似未花费的 output。\footnote{\label{input-selection}在门罗币中，"额外不相关地址"的集合\marginnote{src/wallet/ wallet2.cpp {\tt get\_outs()}}通常从 Gamma\marginnote{{\tt gamma\_picker ::pick()}}[1.2cm] 分布中选择，覆盖历史 output 的范围（仅限 RingCT output，对于 pre-RingCT output 使用三角分布）。此方法使用分箱程序来平滑区块密度的差异。首先计算 RingCT output 过去一年内的平均每区块 output 数量（平均时间 = [\#outputs/\#blocks]*blocktime）。通过 Gamma 分布选择一个 output，然后在其区块内查找并抓取一个随机 output 作为集合成员。\cite{AnalysisOfLinkability}}\footnote{从协议 v12 起，所有交易 input 必须至少有 10 个区块的确认时间（{\tt CRYPTONOTE\_DEFAULT\_TX \_SPENDABLE\_AGE}）。在 v12 之前，核心实现默认使用 10 个区块，但不是强制要求，因此替代钱包可以做出不同选择，有些确实这样做了 \cite{visualizing-monero-vid}。}要签名 input $j$，她将一个 $v$ 大小的集合搅入一个{环 (ring)}中，同时加入她自己第 $j$ 个未花费的一次性地址（放在唯一索引 $\pi$ 处），以及虚假的零承诺，如下：\footnote{在门罗币中，每笔交易的每个环必须具有相同的大小，协议控制每个待花费 output 可以有多少环成员。它随不同的协议版本而改变：v2 2016 年 3 月 $\geq$ 3，v6 2017 年 9 月 $\geq$ 5，v7 2018 年 4 月 $\geq$ 7，v8 2018 年 10 月仅 11。自\marginnote{src/crypto- note\_core/ cryptonote\_ core.cpp {\tt check\_ tx\_inputs\_ ring\_members\_ diff()}} v6 起，每个单独的环不能包含重复成员，但环之间可能存在重复以允许多个 input...

\begin{align*}
    \mathcal{R}_j = \{&\{K^o_{1, j}, (C_{1, j} - C'^a_{\pi, j})\}, \\
    &... \\
    &\{ K^o_{\pi, j}, (C^a_{\pi, j} - C'^a_{\pi, j})\}, \\
    &... \\
    &\{ K^o_{v+1, j}, (C_{v+1, j} - C'^a_{\pi, j})\}\}
\end{align*}

Alice\marginnote{src/ringct/ rctSigs.cpp proveRct- MGSimple()}使用 MLSAG 签名（第 \ref{sec:MLSAG} 节）来签名这个环，她知道 $K^o_{\pi,j}$ 的私钥 $k^o_{\pi,j}$，以及零承诺 $(C^a_{\pi,j}$ - $C'^a_{\pi,j})$ 的私钥 $z_j$。由于零承诺不需要 key image，因此签名的构造中也没有相应的 key image 分量。\footnote{构建和验证签名时排除项 $r_{i,2} \mathcal{H}_p(C_{i, j} - C'^a_{\pi, j}) + c_i \tilde{K}_{z_j}$。}

交易中的每个 input 都使用如上定义的 \(\\mathcal{R}_j\) 等环单独签名，从而将实际花费的 output（$K^o_{\pi,1},...,K^o_{\pi,m}$）隐藏在其他未花费的 output 中。\footnote{单独签名 input 的优势在于，真实 input 和零承诺的集合不必放在相同的索引 $\pi$ 处，而在聚合的情况下则需要如此。这意味着即使一个 input 的来源变得可识别，其他 input 的来源也不会。旧的交易类型 {\tt RCTTypeFull} 使用聚合环签名，将所有环合并为一个，据我们了解正是因此而被弃用。}由于每个环的部分包含一个零承诺，因此使用的伪 output 承诺必须包含与实际花费的 input 相等的金额。这将 input 金额与金额平衡的证明绑定在一起，而不会泄露哪个环成员是真正的 input。

每个 input 签名的消息 $\mathfrak{m}$ 本质上是所有交易数据的 hash，{不包括} MLSAG 签名。\footnote{实际\marginnote{src/ringct/ rctSigs.cpp {\tt get\_pre\_ mlsag\_hash()}}消息是 $\mathfrak{m} = \mathcal{H}(\mathcal{H}(tx\textunderscore prefix),\mathcal{H}(ss),\mathcal{H}(\text{range proofs}))$，其中：\par
$tx\textunderscore prefix = $\{交易时代版本（即 RingCT = 2），input \{环成员 key 偏移，key image\}，output \{一次性地址\}，extra \{交易公钥，支付 ID 或编码支付 ID，杂项。\}\}\par
$ss = $\{交易类型（{\tt RCTTypeBulletproof2} = '4'），交易手续费，input 的伪 output 承诺，ecdhInfo（加密金额），output 承诺\}。\par
参见附录 \ref{appendix:RCTTypeBulletproof2} 关于此术语。}这确保了从交易作者和验证者的角度来看，交易是防篡改的。只产生一个消息，每个 input 的 MLSAG 都签名它。

一次性私钥 $k^o$ 是门罗币交易模型的本质。用 $k^o$ 签名 $\mathfrak{m}$ 证明你是 $C^a$ 中承诺的金额的所有者。验证者可以确信交易作者在花费自己的资金，而甚至不知道在花费哪笔资金、花了多少，或者他们还拥有什么其他资金！


\subsection{避免双花 (Avoiding double-spending)}

一个\marginnote{src/crypto- note\_core/ block- chain.cpp {\tt have\_tx\_ keyimges\_ as\_spent()}} MLSAG 签名（第 \ref{sec:MLSAG} 节）包含私钥 \(k_{\pi, j}\) 的像 \(\\tilde{K}_{j}\)。任何密码签名方案的一个重要性质是它应该是不可伪造的，概率不可忽略。因此，在所有实际效果中，我们可以假设签名的 key image 必须是由合法私钥确定性产生的。

网络只需验证 MLSAG 签名中包含的 key image（对应于 input，计算为 $\tilde{K}^o_{j} = k^o_{\pi,j} \mathcal{H}_p(K^o_{\pi,j})$）没有在其他交易中出现过。\footnote{验证者还必须检查 key image 是否是生成元子群的成员（回顾第 \ref{blsag_note} 节）。}如果出现过，那么我们可以确信我们正在目睹一个重新花费 output $(C^a_{\pi,j}, K_{\pi,j}^o)$ 的企图。


\subsection{空间需求 (Space requirements)}
\label{subsec:space-and-ver-rcttypefull}

\subsubsection*{MLSAG 签名（input）}

从第 \ref{sec:MLSAG} 节我们回顾，在此上下文中的 MLSAG 签名可以表示为

\hfill \(\sigma_j(\mathfrak{m}) = (c_1, r_{1, 1}, r_{1, 2}, ..., r_{v+1, 1}, r_{v+1, 2}) \textrm{ with } \tilde{K}^o_j \) \hfill \phantom{.}

作为 CryptoNote 的遗留问题，值 \(\tilde{K}^o_j\) 不被称为签名的一部分，而是私钥 $k^o_{\pi,j}$ 的{像 (images)}。这些{key image} 通常在交易结构中单独存储，因为它们用于检测双花攻击。

考虑到这一点并假设点压缩（第 \ref{point_compression_section} 节），由于每个环 \(\mathcal{R}_j\) 包含 \((v+1) \cdot 2\) 个密钥，一个 input 签名 $\sigma_j$ 将需要 \( (2(v+1) + 1) \cdot 32  \) 字节。除此之外，key image $\tilde{K}^o_{\pi,j}$ 和伪 output 承诺 $C'^a_{\pi,j}$ 总共留下 $(2(v+1)+3) \cdot 32$ 字节 per input。

对于这个值，我们还需要额外的空间来存储区块链中的环成员偏移量（参见附录 \ref{appendix:RCTTypeBulletproof2}）。这些偏移量被验证者用来在区块链中找到每个 MLSAG 签名的环成员的 output 密钥和承诺，并以可变长度整数存储，因此我们无法精确量化所需的空间。\footnote{参见 \cite{varint-description} 或 \cite{varint-spec} 了解门罗币的 varint 数据类型\marginnote{src/common/ varint.h}。它是一种最多使用 9 字节的整数类型，最多可存储 63 位信息。}\footnote{假设\marginnote{src/crypto- note\_basic/ cryptonote\_ format\_ utils.cpp {\tt absolute\_out- put\_offsets\_ to\_relative()}}区块链包含一个长长的交易 output 列表。我们报告想要在环中使用的 output 的索引。现在，更大的索引需要更多的存储空间。我们只需要从每个环中报告{一个}索引的"绝对"位置，以及其他环成员索引的"相对"位置。例如，对于真实索引 \{7,11,15,20\}，我们只需要报告 \{7,4,4,5\}。验证者可以用 (7+4+4+5 = 20) 计算最后一个索引。环成员在环内按区块链索引升序排列。}\footnote{使用 11 个总成员的环和 10 个 input 的交易将需要 \(((11 \cdot 2 + 3) \cdot 32) \cdot 10 = 8000 \) 字节来存储其 input，大约 110 到 330 字节用于偏移量（有 110 个环成员）。}%src/cryptonote_basic/cryptonote_format_utils.cpp absolute_output_offsets_to_relative

验证\marginnote{src/ringct/ rctSigs.cpp verRctMG- Simple() verRct- Semantics- Simple()}[-1cm]一笔 {\tt RCTTypeBulletproof2} 交易的所有 MLSAG 包括计算 \( (C_{i, j} - C'^a_{\pi, j}) \) 和 \( (\sum_j C'^a_{j} \stackrel{?}{=} \sum_t C^b_{t} + f H)\)，以及验证 key image 在 $G$ 的子群中，即 $l \tilde{K} \stackrel{?}{=} 0$。

\subsubsection*{范围证明（output）}

一个\marginnote{src/ringct/ bullet- proofs.cpp {\tt bullet- proof\_ VERIFY()}}聚合 Bulletproof 范围证明将需要 $(2 \cdot \lceil \textrm{log}_2(64 \cdot p) \rceil + 9) \cdot 32$ 字节的总存储空间。



\newpage
\section{概念总结：门罗币交易 (Concept summary: Monero transactions)}
\label{sec:transaction_summary}

为了总结本章以及前两章，我们以概念清晰的方式呈现交易的主要内容。真实示例可以在附录 \ref{appendix:RCTTypeBulletproof2} 中找到。

\begin{itemize}
    \item \underline{类型 (Type)}：`0' 是 {\tt RCTTypeNull}（用于矿工），`4' 是 {\tt RCTTypeBulletproof2} %see chapter 7, blockchain, about type 0 transactions
    \item \underline{Input}：对于交易作者花费的每个 input $j \in \{1,...,m\}$
    \begin{itemize}
        \item \textbf{环成员偏移量 (Ring member offsets)}：一个"偏移量"列表，指示验证者可以在区块链中找到 input $j$ 的环成员 $i \in \{1,...,v+1\}$（包括真实 input）
        \item \textbf{MLSAG 签名}：$\sigma_j$ 项 $c_1$，以及 $i \in \{1,...,v+1\}$ 的 $r_{i,1}$ \& $r_{i,2}$
        \item \textbf{Key image}：input $j$ 的 key image $\tilde{K}^{o,a}_j$
        \item \textbf{伪 output 承诺 (Pseudo output commitment)}：input $j$ 的 $C'^{a}_j$
    \end{itemize}
    
    \item \underline{Output}：对于发往地址或子地址 $(K^v_t,K^s_t)$ 的每个 output $t \in \{0,...,p-1\}$
    \begin{itemize}
        \item \textbf{一次性地址 (One-time address)}：output $t$ 的 $K^{o,b}_t$
        \item \textbf{Output 承诺}：output $t$ 的 $C^{b}_t$
        \item \textbf{编码金额 (Encoded amount)}：以便 output 所有者可以计算 output $t$ 的 $b_t$
        \begin{itemize}
            \item \textit{金额 (Amount)}：$b_t \oplus_8 \mathcal{H}_n(``amount", \mathcal{H}_n(r K_B^v, t))$
        \end{itemize}
        \item \textbf{范围证明 (Range proof)}：所有 output 金额 $b_t$ 的聚合 Bulletproof
        \begin{itemize}
            \item \textit{证明 (Proof)}：$\Pi_{BP} = (A, S, T_1, T_2, \tau_x, \mu, \mathbb{L}, \mathbb{R}, a, b, t)$
        \end{itemize}
    \end{itemize}
    \item \underline{交易手续费 (Transaction fee)}：以明文传达，乘以 $10^{12}$（即原子单位，见第 \ref{subsec:block-reward} 节），因此 1.0 的手续费将被记录为 1000000000000
    \item \underline{附加数据 (Extra)}：包括交易公钥 $r G$，或者如果至少有一个 output 发往子地址，则包括每个子地址 output $t$ 的 $r_t K^{s,i}_t$ 和每个普通地址 output $t$ 的 $r_t G$，也许还有一个编码的支付 ID（每笔交易最多一个）\footnote{存储在"extra"字段中的信息不会被验证，但它{确实}由 input MLSAG 签名，因此不可能被篡改（概率可忽略不计）。该字段对可以存储多少数据没有限制，只要遵守最大交易权重即可。更多详情请参见 \cite{extra-field-stackexchange}。}
\end{itemize}

我们最后的单 input/单 output 示例交易听起来像这样："我的交易使用交易公钥 $r G$。我将花费集合 $\mathbb{X}$ 中的一个 output（注意它有一个隐藏的金额 $A_X$，在 $C_X$ 中做出承诺）。正在花费的 output 由我拥有（我在 $\mathbb{X}$ 中的一次性地址上做了 MLSAG 签名），并且之前没有被花费过（它的 key image $\tilde{K}$ 尚未出现在区块链中）。我将给它一个伪 output 承诺 $C'_X$。我将创建一个 output $Y$，可由一次性地址 $K^o_Y$ 花费。它有一个隐藏的金额 $A_Y$ 在 $C_Y$ 中做出承诺，为接收者加密，并通过 Bulletproofs 风格的范围证明证明在范围内。我的交易包含交易手续费 $f$。请注意 $C'_X - (C_Y + C_f) = 0$，并且我已经对零承诺 $C'_X - C_X = z G$ 进行了签名，这意味着 input 金额等于 output 金额（$A_X = A'_X = A_Y + f$）。我的 MLSAG 签名了所有交易数据，因此观察者可以确信它没有被篡改。"


\newpage
\subsection{存储需求 (Storage requirements)}

对于 {\tt RCTTypeBulletproof2}，我们需要 $(2(v+1)+2) \cdot m \cdot 32$ 字节的存储空间，聚合 Bulletproof 范围证明需要 $(2 \cdot \lceil \textrm{log}_2(64 \cdot p) \rceil + 9) \cdot 32$ 字节的存储空间。\footnote{交易内容的数量受所谓最大"交易权重"的限制。\marginnote{src/crypto- note\_core/ tx\_pool.cpp {\tt get\_trans- action\_weight \_limit()}}在 Bulletproofs 于协议 v8 中被添加之前（实际上目前当交易只有两个 output 时），交易权重和字节大小是等价的。最大权重为 (0.5*300kB - {\tt CRYPTONOTE\_COINBASE\_BLOB\_RESERVED\_SIZE})，其中 blob 预留（600 字节）旨在为区块中的矿工交易留出空间。在 v8 之前不包含 0.5 倍乘数，300kB 项在早期协议版本中更小（v1 为 20kB，v2 为 60kB，v5 为 300kB）。我们在第 \ref{subsec:dynamic-block-weight} 节中详细说明这些话题。}

杂项需求：
\begin{itemize}
    \setlength\itemsep{\listspace}
    \item Input key image：$m*32$ 字节
    \item 一次性 output 地址：$p*32$ 字节
    \item Output 承诺：$p*32$ 字节
    \item 编码 output 金额：$p*8$ 字节
    \item 交易公钥：通常 32 字节，如果向至少一个子地址发送则为 $p*32$ 字节。
    \item 支付 ID：集成地址为 8 字节。每笔交易不应超过一个。
    \item 交易手续费：以可变长度整数存储，因此 $\leq 9$ 字节
    \item Input 偏移量：以可变长度整数存储，因此每个偏移量 $\leq 9$ 字节，共 $m*(v+1)$ 个环成员
    \item 解锁时间：以可变长度整数存储，因此 $\leq 9$ 字节\footnote{任何交易的作者都可以锁定其 output\marginnote{src/crypto- note\_core/ block- chain.cpp {\tt is\_tx\_ spendtime\_ unlocked()}},使其在指定区块高度之前不可花费（或直到 UNIX 时间戳，届时其 output 可以被纳入区块交易的环成员中）。他只能将所有 output 锁定到相同的区块高度。目前尚不清楚这是否为交易作者提供了有意义的用途（也许是智能合约）。矿工交易有强制的 60 区块锁定时间。从协议 v12 起，普通 output 在默认可花费时间（10 个区块）之前不能被花费，这在功能上等同于强制的最低 10 区块锁定时间。如果一笔交易在第 10 个区块发布，解锁时间为 25，则它可以在第 25 个区块或之后被花费。解锁时间可能是普通交易中最少使用的交易功能。}
    \item "Extra" 标签："extra" 字段中的每条数据（例如交易公钥）以 1 字节的 "tag" 开头，某些数据还有 1+ 字节的 "length"；更多详情请参见附录 \ref{appendix:RCTTypeBulletproof2}
\end{itemize}